\subsection{8. Die fundamentale Frage: was fehlt
wirklich?}\label{die-fundamentale-frage-was-fehlt-wirklich}

\subsubsection{8.1 Ein überprüfbarer
Unterbestimmtheitssatz}\label{ein-uxfcberpruxfcfbarer-unterbestimmtheitssatz}

Betrachte bei festem W dieselbe Familie \(H_{\Delta,g}=\Delta N_b+gX\).
Alle diese Hamiltonoperatoren besitzen die gleiche innere
Quellsymmetrie, dieselbe Ladungserhaltung und denselben kommutierenden
endlichen Clock. Dennoch unterscheiden sie sich dynamisch. Schon für
dieselbe gefüllte Referenz F ist die dimensionslose Momentenkombination

\[\adjustbox{max width=\linewidth}{$\displaystyle 
\frac{\mu_4\mu_2}{\mu_3^2}
=1+1396\left(\frac{g}{\Delta}\right)^2.
$}\]

Bei g geteilt durch Delta gleich 1/20 ist sie 4,49, bei 1/40 dagegen
1,8725. Beide Parameterpunkte liegen im früher abgesicherten
Grundzustandsintervall. Weil die Größe dimensionslos ist, handelt es
sich nicht lediglich um eine Änderung der Maßeinheit. Die gleiche
Grammatik und Symmetrie lässt also unterschiedliche Antworten auf
dieselbe Präparation zu.

\textbf{Damit ist die eindeutige physische Dynamik aus diesen Grunddaten
allein nicht ableitbar.} Das widerlegt nicht die Möglichkeit einer
einfachen zusätzlichen Ursprungsregel. Es identifiziert genau deren
Aufgabe: Sie muss den Prozess samt verfügbarem Operationssatz und
Zustand auswählen, statt nur seine Symmetrie zu benennen. Im fest
gesetzten Prüfvertrag sind g und Delta natürlich definiert; offen ist
ihre physische Herkunft.

\subsubsection{8.2 Der kleinste tragfähige
Forschungsansatz}\label{der-kleinste-tragfuxe4hige-forschungsansatz}

Für eine einzelne Antwortrechnung sollte der relevante Raum nicht von
vornherein der gesamte Fockraum sein. Der zyklische Raum der
tatsächlichen Operationen auf dem angegebenen Startzustand ist der
kleinere Arbeitsraum. Die H-Lanczos-Kette ist eine exakte Darstellung
eines solchen Raums. Mehrere Präparationen, Messungen oder
Compileroperationen können jedoch einen größeren gemeinsamen Prozessraum
verlangen. Ihre Kompatibilität muss dann nachgewiesen werden; getrennte
kleine Modelle dürfen nicht einfach zusammengefügt werden. Eine
vollständige Rekonstruktion des Grundzustands aus dieser Referenz
verlangt außerdem einen Nachweis ihrer nichtverschwindenden Überlappung
mit Omega; die Ritz-Obergrenze allein ersetzt ihn nicht.

Die nächste physikalische Schließung bleibt dieselbe konkrete Bedingung:
Auf einer gemeinsam hergeleiteten Quelle und demselben Zustand muss eine
verfügbare Operation eine geladene Anregung zwischen operational
bestimmten Teilen übertragen. Der Test muss die Mehrzeitantwort von
bloßem Überlapp unterscheiden und ein konsistentes relativistisches
Feldwörterbuch verwenden. Koordinatenwechsel, Gram-Überlapp oder ein
innerer Clock allein erfüllen diese Bedingung nicht.

\subsection{9. Aktueller Stand T1 bis
T8}\label{aktueller-stand-t1-bis-t8}

Die folgende Tabelle ordnet die Beiträge dieser Runde den bisherigen
offenen Beweispflichten zu; sie benennt keine neue Definition der
ursprünglichen Tore.

{\def\LTcaptype{none} % do not increment counter
\begin{longtable}[]{@{}
  >{\raggedright\arraybackslash}p{(\linewidth - 4\tabcolsep) * \real{0.3333}}
  >{\raggedright\arraybackslash}p{(\linewidth - 4\tabcolsep) * \real{0.3333}}
  >{\raggedright\arraybackslash}p{(\linewidth - 4\tabcolsep) * \real{0.3333}}@{}}
\toprule\noalign{}
\begin{minipage}[b]{\linewidth}\raggedright
Tor
\end{minipage} & \begin{minipage}[b]{\linewidth}\raggedright
Beitrag dieser Runde
\end{minipage} & \begin{minipage}[b]{\linewidth}\raggedright
Weiter fehlender entscheidender Nachweis
\end{minipage} \\
\midrule\noalign{}
\endhead
\bottomrule\noalign{}
\endlastfoot
T1 & Exakte bedingte Kontrollalgebra; Auslese- und Eingriffslücke
präzisiert & Tatsächlich erzeugbare Compileroperationen und
Instrumente \\
T2 & Bessere lokale Struktur und Antwortgrenze in der endlichen Bank &
Quellseitiges Half-Charge-Feld mit Energie- und Adjungiertenkontrolle
sowie E8-/Clock-Zuordnung \\
T3 & Kleiner gemeinsamer Rechenraum innerhalb einer Bank & Gemeinsamer
physischer 3+1D-Ursprung \\
T4 & Vollständige innere Bilineare; korrekte Feldkanäle & Vollständiges
chirales Maß mit dynamischen Eichfreiheitsgraden \\
T5 & Exakter lokaler Hamilton-Lanczos-Anfang & Kontrollierter
relativistischer Kontinuums- und Streuungsgrenzfall \\
T6 & Expliziter Nachweis, dass die bisherigen Grunddaten g/Delta nicht
auswählen & Physische Kopplungs-, Skalen- und Spektralzuordnung,
einschließlich offener Neutrinofragen \\
T7 & Keine neue Schließung & Aus derselben Quelle erzeugter dynamischer
Spin-2-Sektor und universelle Kopplung \\
T8 & Unterscheidung Startzustand, Grundzustand und Ausleseinstrument
verschärft & Gemeinsame physische Präparation, Aufzeichnung und
Zustandsauswahl \\
\end{longtable}
}

Alle acht Tore bleiben offen. Die heutigen Ergebnisse schließen
begrenzte mathematische Fragen innerhalb des Programms, nicht diese
physikalischen Gesamtnachweise. Frühere RH-, Faktorisierungs- oder
P-versus-NP-Grenzen werden durch die endlichen Rechnungen ebenfalls
nicht verändert. Der begrenzte RH-Registerabgleich und seine
Aktualitätsgrenze sind in Teil B dokumentiert; ein vollständiger
Neuaudit dieser Arbeitsfronten fand nicht statt.

\subsection{10. Die drei nächsten entscheidenden
Untersuchungen}\label{die-drei-nuxe4chsten-entscheidenden-untersuchungen}

\begin{enumerate}
\def\labelenumi{\arabic{enumi}.}
\tightlist
\item
  \textbf{Den minimalen Ursprungsvertrag auswählen.} Teil B
  klassifiziert die kleinen symmetrieverträglichen Terme und trennt
  echte Erweiterungen von einem bloßen Wechsel der Bosonvariablen. Der
  Compiler muss diese Auswahl und die relativen Parameter begründen.
  Symmetrie allein tut das noch nicht.
\item
  \textbf{Die gemeinsame Schattenkarte am vorhandenen Tensor prüfen.}
  Derselbe Zustand und dieselben Operationen müssen die verschiedenen
  Ausleseansichten tragen. Gesucht sind ergänzende statt nur duplizierte
  Informationen, ein kontrollierter gemeinsamer Kern und Vorhersagen
  ohne nachträgliche Anpassung.
\item
  \textbf{Einen echten Eingriffstransfer und denselben Feldadapter
  konstruieren.} Eine aus der Quelle abgeleitete lokale Operation muss
  später die unbedingte Statistik eines anderen operational bestimmten
  Teils verändern. Selbst zeitabhängige Kreuzkorrelation bei Anfangswert
  null reicht nicht; Teil B enthält ein exaktes Gegenbeispiel. Dazu
  müssen Ladung, Kinetik und relativistische Feldzuordnung passen, bevor
  räumliche Skalierung trägt.
\end{enumerate}

Die korrekte H-Kette mit Restschranke bleibt eine nützliche parallele
Rechenaufgabe. Die zehn bekannten Momente sind feste Gegenprüfungen. Sie
ersetzt aber nicht die jetzt vorrangige Ursprungs- und
Kompositionsfrage.

\subsection{11. Reproduktion und
Veröffentlichung}\label{reproduktion-und-veruxf6ffentlichung}

Das Prüfpaket enthält sechs neue Python-Prüfer, die eingefrorenen
Quellen, die alte und neu wiederholte Kontraktion, Ergebnisdateien,
Quellenmanifest und die vollständige v1.6.6 als historische Basis.
\codeword{replay.py} prüft alle Pins, beide Python-Modi und die
C++-Kontraktion. Die ausgelieferte ZIP wird nach dem Entpacken erneut
geprüft.

Die vollständige Hauptfassung, das kurze Änderungsdokument und die
einfache Erklärung werden gemeinsam versioniert direkt in Documents
abgelegt, jeweils als Markdown und PDF. Die PDF-Ausgaben werden
zusätzlich auf Darstellung, Textabdeckung und Seitenaufbau geprüft.
Frühere Fassungen bleiben erhalten. Es erfolgt keine Änderung der
zentralen TOE-Abnahmemarker, kein Commit, kein Push und keine externe
Web-Veröffentlichung.

Der Statusleitfaden beeinflusst die Trennung zwischen mathematischer
Struktur, angenommenen Ressourcen und nachgewiesener Fähigkeit. Der
Debugging-Leitfaden führte zur ausdrücklichen Sicherung
reell-ganzzahliger Tensorwerte vor der Typumwandlung; keine
Imaginärteile werden still verworfen.

\subsubsection{Schluss in einem Satz}\label{schluss-in-einem-satz}

\textbf{Die einfache Regel ist vorhanden; ihr erster verzweigter
Zustandsbereich und die richtige kleine Rechenkette sind jetzt genauer
gelöst. Was noch fehlt, ist die aus dem Ursprung folgende Auswahl und
Ausführung dieses Prozesses als dieselbe physikalische Welt.}

\begin{center}\rule{0.5\linewidth}{0.5pt}\end{center}

\section{Neue Vogelperspektive: derselbe Prozess in verschiedenen
Ansichten}\label{neue-vogelperspektive-derselbe-prozess-in-verschiedenen-ansichten}

\subsection{B1. Was sich durch die drei weiteren Texte
ändert}\label{b1.-was-sich-durch-die-drei-weiteren-texte-uxe4ndert}

Die nachgereichten Texte werden als Vorschläge geprüft, nicht als
bewiesene Ergebnisse oder neue Arbeitsanweisungen übernommen. Die
stärkste gemeinsame Idee trägt: \textbf{Innere Identität,
Zusammensetzung, Ort, Zeit und Auslese dürfen nicht stillschweigend
dasselbe Label erhalten.} Eine immer genauere Rechnung innerhalb einer
Bank ersetzt diese Verbindungen nicht. Die voranstehenden lokalen
Ergebnisse bleiben gültig; die fundamentale Priorität verschiebt sich
auf die gemeinsame Operations- und Kompositionsregel.

{\def\LTcaptype{none} % do not increment counter
\begin{longtable}[]{@{}
  >{\raggedright\arraybackslash}p{(\linewidth - 2\tabcolsep) * \real{0.5000}}
  >{\raggedright\arraybackslash}p{(\linewidth - 2\tabcolsep) * \real{0.5000}}@{}}
\toprule\noalign{}
\begin{minipage}[b]{\linewidth}\raggedright
Vorschlag
\end{minipage} & \begin{minipage}[b]{\linewidth}\raggedright
Prüfung und Konsequenz
\end{minipage} \\
\midrule\noalign{}
\endhead
\bottomrule\noalign{}
\endlastfoot
Ort in Multiplizitäten statt in den 64 inneren Labels suchen & Sinnvolle
Möglichkeit, aber Multiplizität beschreibt bereits chemische Umwandlung
ohne Raum \\
Der große Kommutant könnte verborgenen Raum enthalten & Hamilton- und
Symmetriekommutant sind verschieden; die Größe allein trägt diesen
Schluss nicht \\
Den ganzen Z4-Ursprung statt nur den ladungserhaltenden Ausschnitt
betrachten & Konkrete zusätzliche Terme existieren am selben W; ihre
Verfügbarkeit ist noch nicht hergeleitet \\
Ladung +3 und Ladung -1 sind modulo vier gleich & Ein zusätzlicher Term
kann sie dynamisch verbinden; die Operatoren sind nicht identisch \\
Zeitabhängige Kreuzantwort zeigt Transport & Im Allgemeinen falsch;
unten steht ein exaktes Gegenbeispiel sogar bei Anfangskorrelation
null \\
Mehrere Schatten könnten einen gemeinsamen Ursprung bestimmen & Möglich
bei gemeinsamem Vertrag und ausreichender Unterscheidbarkeit;
Übereinstimmung allein genügt nicht \\
Holografie könnte die fehlende Verbindung sein & Ein nativer
isometrischer Kodierungsbaustein ist vorhanden; Raumgrenze, Schutz und
gemeinsame Raumzeit fehlen \\
\end{longtable}
}

Die zusätzlich eingefrorenen Texte heißen im Prüfpaket
\codeword{multiplicity\_z4.txt}, \codeword{fundamental\_reduction.txt}
und \codeword{runtime\_synthesis.txt}. Sie werden im Hauptdokument
vollständig angehängt. Die anschließende Schatten- und Holografiefrage
des Nutzers wurde durch eigene Rechnungen weiterverfolgt.

\subsection{B2. Noethers Frage: Welche Unterschiede kann die Dynamik
überhaupt
verändern?}\label{b2.-noethers-frage-welche-unterschiede-kann-die-dynamik-uxfcberhaupt-veruxe4ndern}

Bei ungebrochener innerer Gruppe G zerfällt ein Zustandsraum als

\[\adjustbox{max width=\linewidth}{$\displaystyle 
\mathcal H=\bigoplus_\lambda V_\lambda\otimes M_\lambda,
\qquad H=\bigoplus_\lambda I_{V_\lambda}\otimes h_\lambda,
\qquad \operatorname{End}_G(\mathcal H)
=\bigoplus_\lambda I_{V_\lambda}\otimes\operatorname{End}(M_\lambda).
$}\]

G-invariante Entwicklung darf die Kopien eines Typs verändern, nicht die
inneren Komponenten dieses Typs willkürlich auseinanderziehen. Das macht
Multiplizitäten zu einem möglichen Träger weiterer Dynamik. Es macht sie
noch nicht zu Orten. Schon die drei hellen N=3-Typen enthalten zwei
Kopien: einmal drei Fermionen, einmal einen Fermion plus einen Boson.
Ihre jeweilige Hamiltonmatrix lautet

\[\adjustbox{max width=\linewidth}{$\displaystyle 
h_\lambda=\begin{pmatrix}0&g\sqrt\lambda\\g\sqrt\lambda&\Delta\end{pmatrix},
\qquad \lambda=7,10,12.
$}\]

Am Prüfpunkt g/Delta gleich 1/20 sind die maximalen
Umwandlungswahrscheinlichkeiten 7/107, 1/11 und 3/28. Das ist echte
symmetrieverträgliche Dynamik, zunächst aber eine Veränderung der
Zusammensetzung und kein Ortswechsel.

Wichtig ist die Richtung des Kommutanten: Im N=3-Sektor hat der
Kommutant der vollen inneren Gruppe die Dimension 16. Die viel größere
Zahl 1444233216 betrifft den kleineren Kontrollsatz aus X und Bosonzahl.
Auf dem isolierten irreduziblen 64er-Pol ist der G-Kommutant sogar nur
skalar, während der Kommutant der dort skalaren Energie alle
64-mal-64-Matrizen enthält. Die große Energiedegeneration ist kein
Nachweis vieler Orte.

Für einen G-Singulett-Grundzustand und die ursprünglichen irreduziblen
Fermionoperatoren folgt aus Schurs Lemma

\[\adjustbox{max width=\linewidth}{$\displaystyle 
C_{rs}(t)=\delta_{rs}\,c(t).
$}\]

Feste lineare Umbenennungen liefern lediglich einen Gramfaktor mal
derselben Zeitfunktion. Auch eine perfekte Kenntnis aller lokalen
Spektrallinien würde aus den 64 inneren Komponenten deshalb noch keine
räumliche Anordnung machen. Ein stationärer Grundzustand kann
zeitabhängige Korrelationsfunktionen haben; Nichtstationarität ist keine
notwendige Voraussetzung dafür. Physische Uhrenablesung und Zeitrichtung
sind weitere Fragen.

\subsection{B3. Korrelation ist noch kein Eingriff: ein exakt gelöstes
Gegenbeispiel}\label{b3.-korrelation-ist-noch-kein-eingriff-ein-exakt-geluxf6stes-gegenbeispiel}

Zwei nicht miteinander wechselwirkende Qubits genügen:

\[\adjustbox{max width=\linewidth}{$\displaystyle 
H=\tfrac12(Z_A+Z_B),\qquad
|\Omega\rangle=(|01\rangle+|10\rangle)/\sqrt2,
\qquad A=X_A,\quad B=Y_B.
$}\]

Der Zustand ist stationär. Trotzdem ist

\[\adjustbox{max width=\linewidth}{$\displaystyle 
\langle\Omega|B e^{-itH}A|\Omega\rangle=\sin t.
$}\]

Die Kreuzkorrelation startet bei null und verändert sich mit der Zeit.
Es gibt dennoch keinerlei Signal von A nach B: Alle A-lokalen
Operationen kommutieren mit jedem zeitentwickelten B-lokalen
Observablen. Für jedes spurtreue lokale Quanteninstrument bleibt daher
die unbedingte B-Statistik unverändert. Eine Konditionierung auf ein
Messergebnis in A wäre etwas anderes.

Der belastbare Übertragungstest ist stattdessen

\[\adjustbox{max width=\linewidth}{$\displaystyle 
\delta_{A\to B}(t)=
\operatorname{tr}\!\left[B U_t\mathcal E_A(\rho)U_t^\dagger\right]
-\operatorname{tr}\!\left[B U_t\rho U_t^\dagger\right].
$}\]

A und B müssen zuerst als unabhängig adressierbare Teile definiert sein;
das Instrument muss verfügbar sein, ohne durch seine Definition schon B
anzusteuern. Erfolg bedeutet eine nichtverschwindende unbedingte
Änderung nach dem A-Eingriff. Die lineare Antwort misst entsprechend
einen retardierten Kommutator, nicht bloß eine Zweipunktkorrelation.
Lieb-Robinson-Abschätzungen können anschließend die Ausbreitung in einem
bereits lokal gekoppelten Modell begrenzen; sie erzeugen diese Lokalität
oder drei Raumdimensionen nicht von selbst. Siehe die Primärarbeit von
\href{https://arxiv.org/abs/quant-ph/0603121}{Bravyi, Hastings und
Verstraete}.

Ein weiterer exakter Befund verhindert ein zu starkes Paritäts-No-go:
Der native Graph der 480 Fermionpaare verbindet alle 64 Moden. Eine
Vorzeichenumkehr auf einer Fermionteilmenge, bei unveränderten Bosonen,
erhält jede Kante genau dann, wenn sie trivial oder die globale
Fermionparität ist. Die binäre Inzidenzmatrix hat Rang 63. Zusätzliche
unabhängige Bankparitäten entstehen beim Kopieren von Banken; sie sind
nicht als solche ein Hindernis innerhalb der einen ursprünglichen Bank.
Werden auch Bosonen umgezeichnet, sind andere gemeinsame Gradierungen
möglich.

\subsection{B4. Die einfachste fehlende Auswahl: Warum genau dieser
Hamiltonoperator?}\label{b4.-die-einfachste-fehlende-auswahl-warum-genau-dieser-hamiltonoperator}

Die Z4-Idee besitzt einen präzisen Kern. Ein ganzzahliger, auf jedem
Grad konstanter Ladungslift mit neutralem Grad null müsste bei den
entsprechenden nichtverschwindenden Klammern erfüllen

\[\adjustbox{max width=\linewidth}{$\displaystyle 
2q_1=q_2,\qquad q_1+q_3=0,\qquad 2q_3=q_2.
$}\]

Die Koeffizientenmatrix hat Determinante 4 und Rang drei. Über den
ganzen Zahlen bleibt nur q1=q2=q3=0; modulo vier gilt dagegen die
gewohnte Belegung 1, 2, 3. Das verbietet nicht andere kontinuierliche
Cartanladungen und beweist nicht, dass der aktuelle Operator-Hamiltonian
schon alle E8-Klammern realisiert.

Eine wörtliche Identifikation der Grad-eins-Labels mit primitiven CAR-
Vernichtern scheitert bereits an einem kleinen Beispiel: Die native
W-Spalte des Paares (0,1) ist null, aber \([f_0,f_1]=2f_0f_1\ne0\). Die
innere Lieklammer und die gewöhnliche Operatoralgebra der Fermionfelder
dürfen nicht gleichgesetzt werden. Eine andere, zusammengesetzte
E8-Realisierung wird dadurch nicht ausgeschlossen.

